% Generated draft: requires mathematical and source review.
\documentclass[11pt]{article}
\usepackage[margin=25mm]{geometry}
\usepackage{fontspec}
\setmainfont{Latin Modern Roman}
\usepackage{amsmath,amssymb,mathtools,unicode-math}
\setmathfont{Latin Modern Math}
\usepackage{enumitem,needspace,xurl,hyperref}
\setlength{\parindent}{0pt}
\setlength{\parskip}{5pt}
\setlength{\emergencystretch}{3em}
\newcommand{\N}{\mathbb{N}}
\newcommand{\Z}{\mathbb{Z}}
\newcommand{\Q}{\mathbb{Q}}
\newcommand{\R}{\mathbb{R}}
\newcommand{\C}{\mathbb{C}}
\newcommand{\F}{\mathbb{F}}
\newcommand{\A}{\mathbb{A}}
\newcommand{\PP}{\mathbb{P}}
\newcommand{\E}{\mathbb{E}}
\newcommand{\eps}{\varepsilon}
\newcommand{\dd}{\,\mathrm{d}}
\begin{document}
\section*{UnsolvedMath ID: 30002009}
\textbf{Source problem code:} catalogue entry
\subsection{Definitions and mathematical statement}
Let \(G\) be a complex reductive algebraic group, let \(T\subseteq G\) be a maximal torus, and let \(W=N_G(T)/T\) be its Weyl group. Let \(k\) be an arbitrary field, with no restriction on its characteristic, and write \(k[W]\) for the group algebra.

Under the usual topological interpretation, put \(X=(G/T)(\C)\), with its complex analytic topology, and let \(k_X\) denote the constant sheaf with value \(k\). The right action of \(W\) on \(X\) is given by
\[
(gT)\cdot(nT)=gnT,\qquad n\in N_G(T).
\]
Pullback along this right action induces a left \(k[W]\)-module structure on derived global sections. Thus the object to be described is
\[
C_{G,k}:=R\Gamma(X,k_X)\in D(k[W]),
\]
where \(D(k[W])\) denotes the derived category of left \(k[W]\)-modules. A complex is called perfect if it is quasi-isomorphic to a bounded complex of finitely generated projective \(k[W]\)-modules.

\textbf{Problem.} For every such \(G\) and \(T\), and every field \(k\), describe the perfect complex \(C_{G,k}\) up to isomorphism in \(D(k[W])\). The target is the derived object, including its \(W\)-action, rather than merely its cohomology groups.

\textbf{Unresolved source qualifications.} The supplied record does not explicitly specify the cohomology theory or topology, the convention for the induced \(W\)-action, or whether ``reductive group'' requires connectedness. The construction above makes the usual analytic interpretation explicit; these conventions require confirmation against the source. The record also does not specify the desired form of a general description. Perfectness is a premise reported by the supplied source record, not a result proved in this draft.
\subsection{Short English statement}
For a complex reductive group \(G\), a maximal torus \(T\), and any field \(k\), describe the derived cohomology of \(G/T\) together with its Weyl group action, as a perfect complex of \(k[W]\)-modules. The analytic cohomology convention and the connectedness convention require source confirmation.
\subsection{Sources}
\begin{itemize}
\item Supplied local UnsolvedMath registry record, permanent ID \(30002009\), source problem code \texttt{catalogue entry}, in the collection \emph{Oberwolfach Reports Open Problems}. The record attributes the problem to the 2012 workshop report; no individual proposer is supplied.
\item \emph{Enveloping Algebras and Geometric Representation Theory} (2012), Oberwolfach Reports, DOI \texttt{10.4171/owr/2012/13}, pp.~733--809. Supplied source link and literature-source label ``Oberwolfach Report 13 2012'': \url{https://doi.org/10.4171/owr/2012/13}.
\item Supplied literature source: \emph{Modular Perverse Sheaves on Flag Varieties II}, \url{https://arxiv.org/abs/1401.7256}.
\item Supplied literature assessment dated exactly \texttt{2026-08-21}: the record labels the problem open and reports that the complex is identified as a graded regular representation over \(\C\), is perfect over arbitrary fields, and has no established general description over an arbitrary field. This is supplied dated curation derived from the OWR record, not independently verified here; the status and cited mathematical claims require independent verification.
\end{itemize}
\end{document}
